\section{Modello Architetturale}

A.L.I.C.E. e' una web application full-stack basata su \textbf{Next.js App Router} e \textbf{Supabase}. Il progetto non usa un application server separato: Next.js gestisce pagine, componenti server, componenti client, Server Actions, route API e proxy; Supabase gestisce autenticazione, database PostgreSQL, API PostgREST e Row Level Security.

\subsection{Architettura hardware}

L'architettura hardware e' volutamente leggera. Gli utenti accedono tramite browser, mentre il codice applicativo viene eseguito su un ambiente serverless compatibile con Next.js. I dati persistenti e le sessioni sono gestiti da Supabase; il suggerimento AI viene generato tramite Groq API.

\begin{table}[H]
\centering
\renewcommand{\arraystretch}{1.25}
\begin{tabularx}{\textwidth}{|l|l|X|}
\hline
\textbf{Componente} & \textbf{Servizio} & \textbf{Responsabilita'} \\
\hline
Client & Browser web & Visualizzazione interfaccia e invio azioni utente. \\
\hline
Web app & Next.js serverless & Rendering, routing, Server Actions, proxy e route API. \\
\hline
Backend gestito & Supabase & Auth, PostgreSQL, PostgREST e policy RLS. \\
\hline
Servizio AI & Groq API & Generazione del suggerimento testuale per l'operatore. \\
\hline
\end{tabularx}
\caption{Componenti hardware e servizi esterni}
\label{tab:arch-hw}
\end{table}

\subsection{Modello architetturale}

Il modello architetturale non e' un MVC classico. La separazione delle responsabilita' e' ottenuta tramite il modello di Next.js App Router:

\begin{itemize}
    \item le \textbf{pagine} definiscono le viste raggiungibili dall'utente;
    \item i \textbf{Server Components} caricano dati e preparano la UI lato server;
    \item i \textbf{Client Components} gestiscono interazioni complesse nel browser;
    \item le \textbf{Server Actions} gestiscono mutazioni come login, registrazione, farmaci ed esercizi;
    \item il \textbf{proxy} gestisce sessione e redirect delle aree private;
    \item Supabase rappresenta il livello dati e applica le policy di visibilita'.
\end{itemize}

Questa architettura riduce il numero di file necessari rispetto a un backend separato e mantiene la logica vicino alle pagine che la utilizzano.

\subsection{Architettura software}

\begin{figure}[H]
\centering
\resizebox{0.98\textwidth}{!}{
\begin{tikzpicture}[
    layer/.style={rectangle, draw=aliceblue, fill=alicelight, thick, rounded corners=4pt, minimum width=11cm, minimum height=1.35cm, align=center},
    comp/.style={rectangle, draw=gray, fill=white, thick, rounded corners=3pt, minimum width=2.7cm, minimum height=0.75cm, align=center, font=\scriptsize},
    ext/.style={rectangle, draw=aliceorange, fill=aliceorange!10, thick, rounded corners=3pt, minimum width=4.2cm, minimum height=0.85cm, align=center, font=\scriptsize\bfseries},
    line/.style={thick, color=aliceblue, -{Stealth[length=2.5mm]}},
]

\node[layer] (browser) at (0, 0) {Browser: Operatore e Paziente};

\node[layer] (next) at (0, -2) {};
\node[comp] at (-3.6, -2) {App Router\\pagine e layout};
\node[comp] at (-1.2, -2) {Server\\Components};
\node[comp] at (1.2, -2) {Client\\Components};
\node[comp] at (3.6, -2) {\texttt{proxy.ts}\\redirect};

\node[layer] (server) at (0, -4) {};
\node[comp] at (-2.8, -4) {Server Actions\\CRUD e auth};
\node[comp] at (0, -4) {Route Handler\\AI};
\node[comp] at (2.8, -4) {Helper\\\texttt{src/lib}};

\node[ext] (supabase) at (-3.0, -6) {Supabase\\Auth + PostgreSQL + RLS};
\node[ext] (groq) at (3.0, -6) {Groq API\\Suggerimento operatore};

\draw[line] (browser) -- (next);
\draw[line] (next) -- (server);
\draw[line] (server) -- (supabase);
\draw[line] (server) -- (groq);
\end{tikzpicture}
}
\caption{Architettura software a livelli del sistema A.L.I.C.E.}
\label{fig:architettura-software}
\end{figure}

\subsection{Responsabilita' dei livelli}

\begin{table}[H]
\centering
\renewcommand{\arraystretch}{1.25}
\begin{tabularx}{\textwidth}{|l|X|}
\hline
\textbf{Livello} & \textbf{Responsabilita'} \\
\hline
Presentazione & Componenti React, form, dashboard, card e pagine semplificate del paziente. \\
\hline
Navigazione & App Router per le route e \texttt{proxy.ts} per redirect e sessione. \\
\hline
Logica applicativa & Server Actions, pagine server e helper in \texttt{src/lib}. \\
\hline
Dati & Tabelle Supabase, funzioni SQL, trigger e policy RLS. \\
\hline
Integrazione AI & Route \texttt{/api/ai/suggerimento-operatore}, chiamata server-side a Groq. \\
\hline
\end{tabularx}
\caption{Responsabilita' dei livelli software}
\label{tab:livelli-software}
\end{table}

\subsection{Stack tecnologico}

\begin{table}[H]
\centering
\renewcommand{\arraystretch}{1.25}
\begin{tabularx}{\textwidth}{|l|l|X|}
\hline
\textbf{Tecnologia} & \textbf{Versione} & \textbf{Uso nel progetto} \\
\hline
\texttt{next} & 16.1.6 & Framework full-stack, App Router, build e proxy. \\
\hline
\texttt{react} / \texttt{react-dom} & 19.2.3 & Componenti dell'interfaccia utente. \\
\hline
\texttt{@supabase/ssr} & 0.8.0 & Client Supabase compatibile con server, client e proxy. \\
\hline
\texttt{typescript} & 5.x & Tipizzazione del codice applicativo. \\
\hline
CSS Modules / CSS globale & n.d. & Stile delle pagine senza librerie UI esterne. \\
\hline
Groq API & via \texttt{fetch} & Suggerimento AI lato server. \\
\hline
\end{tabularx}
\caption{Stack tecnologico realmente presente nel progetto}
\label{tab:stack-tecnologico}
\end{table}

\subsection{DBMS, browser e sistemi supportati}

Il DBMS e' PostgreSQL tramite Supabase. Il codice applicativo usa il client Supabase per leggere e scrivere dati, mentre le regole RLS nel database impediscono l'accesso a dati non autorizzati. L'applicazione e' fruibile da browser moderni su Windows, macOS, Linux, Android e iPadOS. Per il paziente e' consigliato un tablet o un PC con schermo ampio; per l'operatore e' sufficiente un PC, laptop o tablet.

\newpage
